\section{Exam Strategy and Mixed Task Recipes}

\subsection{Ctrl+F keywords: multiple choice, statement validity, units, limiting case, answer check, traps, Danish keywords}
Use for cross-topic exam tactics that recur in marked multiple-choice and calculation problems.

\subsection{Symbols and notation}
\referencetable{tab:strategy:symbols}{Symbols used in cross-topic checking strategies.}
\begin{tabular}{@{}>{$}l<{$} p{10.8cm}@{}}
\toprule
\text{Symbol} & \text{Meaning and usage}\\
\midrule
\mathrm{pH},\ [\ce{H3O+}] & pH and hydronium concentration used in logarithmic answer checks.\\
\Delta G^\circ,\ K & Standard Gibbs energy change and equilibrium constant.\\
E,\ E^\circ,\ n,\ Q & Cell potential, standard potential, transferred-electron amount, and reaction quotient in Nernst checks.\\
n_{\mathrm{gas}},\ p,\ V,\ R,\ T & Gas amount, pressure, volume, gas constant, and absolute temperature.\\
\bottomrule
\end{tabular}

\subsection{Exam recipe: multiple-choice statement elimination}

\subsubsection{One false clause makes an option false}
\begin{equation}\label{eq:strategy:all-clauses-true}
\text{valid answer option}=\text{all required clauses true}
\end{equation}
For statement options, test each clause independently: sign, unit, condition, and exception. Eliminate options with one wrong condition even if the formula name looks right.

\subsubsection{Dimensional and unit check}
\begin{equation}\label{eq:strategy:unit-check}
\text{final unit must match requested quantity}
\end{equation}
Use units as algebraic factors. Common exam errors are Celsius in gas/Arrhenius equations, entropy in \(\mathrm J\) with enthalpy in \(\mathrm{kJ}\), and cell length units in density.

\subsubsection{Logarithm checks for pH, \(K\), and Nernst}
\begin{equation}\label{eq:strategy:log-checks}
\mathrm{pH}=-\log[\ce{H3O+}], \qquad
\Delta G^\circ=-RT\ln K, \qquad
E=E^\circ-\frac{0.0592}{n}\log Q
\end{equation}
A factor \(10\) in concentration changes pH or Nernst terms by one logarithmic unit. Check whether \(\ln\) or \(\log\) belongs with the numerical constant.

\textbf{Use when:} the question asks which statement or numerical option is correct, or several options differ only in sign, unit, condition, or direction.

\textbf{Given / Find:} one or more candidate answers; determine which option satisfies every required condition.

\textbf{Required references:} Equation~\eqref{eq:strategy:all-clauses-true}, Equation~\eqref{eq:strategy:unit-check}, Equation~\eqref{eq:strategy:log-checks}, and the numbered topic equations or tables named by each option.

\textbf{Steps:}
\begin{enumerate}
    \item Split each answer option into separate claims, for example formula choice, sign, unit, phase inclusion, direction, or stated assumption.
    \item For each formula claim, locate the matching numbered topic relation and substitute or compare exactly as that relation states; do not accept an option because its topic name sounds plausible.
    \item For each numerical claim, carry units through the calculation and eliminate the option if its final unit fails Equation~\eqref{eq:strategy:unit-check}.
    \item For pH, equilibrium, or electrochemical logarithm claims, verify logarithm base and sign using Equation~\eqref{eq:strategy:log-checks}.
    \item Eliminate an option immediately when any single claim is false, according to Equation~\eqref{eq:strategy:all-clauses-true}; retain only options whose claims all pass.
\end{enumerate}

\textbf{Checks:} Re-read qualifiers such as ``standard'', ``at equilibrium'', ``spontaneous'', ``largest'', or ``lowest'' before submitting the retained option.

\textbf{Common traps:} checking only the first clause of a compound option; accepting a correct formula with wrong units or wrong logarithm sign; overlooking a stated condition.

\subsection{Recurring mixed recipes}

\subsubsection{Ideal gas stoichiometry}
\begin{equation}\label{eq:strategy:gas-moles}
n_{\mathrm{gas}}=\frac{pV}{RT}
\end{equation}
Convert gas data to moles before applying reaction coefficients. This links gas-law questions with limiting-reactant calculations.

\subsubsection{Percent change in hydronium concentration from pH}
\begin{equation}\label{eq:strategy:ph-percent-change}
\frac{[\ce{H3O+}]_2}{[\ce{H3O+}]_1}=10^{\mathrm{pH}_1-\mathrm{pH}_2}, \qquad
\%\text{ change}=(\text{ratio}-1)100\%
\end{equation}
Use for ocean-acidification or pH-change wording. Lower pH means larger hydronium concentration.

\subsubsection{Recognition keywords, English and Danish}
\referencetable{tab:strategy:keywords}{Recognition keywords and the topic recipe they identify.}
\begin{center}
\begin{tabular}{c c}
\toprule
Exam wording & Go to\\
\midrule
bundfald, precipitate, sparingly soluble & \(K_{sp}, Q_{sp}\)\\
puffer, buffer, Henderson & acid/base buffer\\
galvanisk, electrode, concentration cell & Nernst/electrochemistry\\
planar, bond angle, pyramidal & Lewis/VSEPR\\
limiting, excess, yield, udbytte & stoichiometry\\
vapor pressure, boiling point & intermolecular forces\\
rate law, half-life, Arrhenius & kinetics\\
\bottomrule
\end{tabular}
\end{center}
Use the wording to jump to the relevant recipe before doing algebra.
